\addcontentsline{toc}{subsection}{necessary and sufficient stuff}
 \begin{thm}[Hinr. und notw. Bedingungen für additive difference Metrik]\label{hinr-notw-beds}~\\
   Suppose $\langle A, \succsim \rangle$ is an additive-difference proximity structure such that $f$ is the additive-difference representation from \ref{add-diff-prox-struct-repr}, i.e.,
\[
f(a, b) = \sum_{i=1}^{n} F_i(|\varphi_i(a) - \varphi_i(b)|).
\]
Suppose further that the domain of each $F_i$ is a real interval and $F_i(0) = 0$, $i = 1, \ldots, n$ (through transformation, since $F_i$ are interval scales).

For $\langle A, \succsim \rangle$ to be compatible with a metric, and additionally be additive on the coordinate axes,
\begin{enumerate}
  \item it is necessary that there exists a continuous, convex function $F$, such that $F_i(x) = F(t_i x), t_i > 0,$ for all $x$ in the domain of $F_i$, $i = 1, \ldots, n$; and
  \item sufficient that, in addition (to the conditions in 1.), the function $\log F(e^x)$ is convex.
\end{enumerate}
\end{thm}

\begin{rem}[Additiv auf den Koordinatenachsen]~\\
Additivität auf den Koordinatenachsen bedeutet hier: Wenn $a, b, c \in A$ sich nur auf einem Faktor $A_i$ unterscheiden, dann gilt die $\triangle$-Ungleichung der Metrik mit Gleichheit. In Mathe Sprache heißt das:
$$ a_j = b_j = c_j \forall j \neq i \text{ und } (a, c) \succsim (a, b), (b,c)  \Rightarrow \delta(a, b) + \delta(b, c) = \delta(a, c) $$
\end{rem}

\begin{cor}[Form von $\delta$]~\\
  Seien die Voraussetzungen wie in Satz \ref{hinr-notw-beds}.
  Dann wird die Metrik $\delta$ dargestellt durch
  \begin{equation}
    \delta(a,b) = F^{-1} \Big \{ \sum_{i=1}^{n} F(| \varphi_i(a_i) - \varphi_i(b_i) |) \Big \} \label{cor-delta-form} 
\end{equation}
  wobei $F$ stetig, konvex und steigend ist.

\end{cor}

\begin{proof}[Beweis von 1.]~\\

Nach den Annahmen des Satzes existiert eine Metrik $\delta$ auf $A$ und eine steigende Funktion $G$, so dass für alle $a, b \in A$,

$$
\delta(a, b) = G \left( \sum_{i=1}^{n} F_i \left( |\varphi_i(a_i) - \varphi_i(b_i)| \right) \right).
$$

Angenommen, $a$, $b$ und $c$ unterscheiden sich nur im $i$-ten Faktor und $a|b|c$ (Also wir können Additivität auf der $i$-ten Koordinatenachse für $a, b, c$ benutzen). Daher gilt

\[
G \left( F_i \left( |\varphi_i(a_i) - \varphi_i(b_i)| \right) \right) + G \left( F_i \left( |\varphi_i(b_i) - \varphi_i(c_i)| \right) \right) = G \left( F_i \left( |\varphi_i(a_i) - \varphi_i(c_i)| \right) \right).
\]

Definiere $|\varphi_i(a_i) - \varphi_i(b_i)| = x$ und $|\varphi_i(b_i) - \varphi_i(c_i)| = y$, was ergibt

\[
G[F_i(x)] + G[F_i(y)] = G[F_i(x + y)].
\]

Definiere $H_i(x) = G[F_i(x)]$. Damit reduziert sich die obige Gleichung auf $H_i(x) + H_i(y) = H_i(x + y)$ für alle $x, y$ im Definitionsbereich von $F_i, 1 \leq i \leq n$ und $H_i$ muss linear sein. Da die Werte positiv sind, ist die einzige monotone Lösung dieser Gleichung $H_i(x) = t_i x$, $t_i > 0$. Da $G$ steigend ist, existiert $G^{-1}$ und es gilt

\[
G^{-1}[H_i(x)] = G^{-1}[t_i x] = F_i(x)
\]

wenn $F_i$ definiert ist. Da $G^{-1}$ unabhängig von $i$ ist, können die $F_i$ sich nur im Definitionsbereich bzw. im Wertebereich unterscheiden. ObdA können wir annehmen, dass der Wertebereich von $F_n$ am größten ist. Da der Bereich von $F_n$ den Bereich von $F_i$ umfasst, können wir $F = F_n$ setzen und erhalten damit $F_i(x) = F_n(t_i x) = F(t_i x)$ für $t_i > 0$ und für jedes $x$ im Definitionsbereich von $F_i$.
Stetigkeit und Konvexität werden in den nächsten Lemmata bewiesen.
\end{proof}

\todo{sagen dass conditions wie oben sind}
\begin{lem}\label{lemma-ungl}
  $F(\alpha + \beta + \gamma) + F(\gamma) \geq F(\alpha + \gamma) + F(\beta + \gamma)$.
\end{lem}

\begin{proof}

In der Dreiecksungleichung,

\[
F^{-1} \left[ \sum_{i=1}^{n} F(\alpha_i) \right] + F^{-1} \left[ \sum_{i=1}^{n} F(\beta_i) \right] \geq F^{-1} \left[ \sum_{i=1}^{n} F(\alpha_i + \beta_i) \right],
\]

seien $\alpha_1 = \gamma, \alpha_2 = F^{-1}(F(\alpha + \gamma) - F(\gamma)), \alpha_j = 0$ für $3 \leq j \leq n$ und  $\beta_1 = \beta, \beta_j = 0$ für $2 \leq j \leq n$. Damit gilt,

\begin{align*}
  F(\alpha + \beta + \gamma) &= F \left[ F^{-1} \left[ \sum_{i=1}^{n} F(\alpha_i) \right] + F^{-1} \left[ \sum_{i=1}^{n} F(\beta_i) \right] \right] \\
                             &\geq \sum_{i=1}^{n} F(\alpha_i + \beta_i) \quad (\text{durch die Dreiecksungleichung})\\
                             &= F(\beta + \gamma) + F(\alpha + \gamma) - F(\gamma).
\end{align*}
\end{proof}

\begin{lem}
  $F$ is continuous.\label{lemma-stetig}
\end{lem}

\begin{proof}

Da $F$ steigend ist, sind alle Unstetigkeiten Sprungstellen. Aus Analysis wissen wir, dass reelle monotone Funktionen höchstens abzählbar viele Sprungstellen (Unstetigkeitesstellen erster Art) haben. Dies führen wir auf einen Widerspruch.\\
Angenommen, $F(\alpha)$ ist die untere Grenze einer Sprungstelle, dann existiert $\epsilon > 0$, so dass für alle $\delta > 0$, $F(\alpha + \delta) - F(\alpha) > \epsilon$. Nach Lemma 4 gilt jedoch für jedes $\beta > 0$,

\[
F(\alpha + \beta + \delta) - F(\alpha + \beta) \geq F(\alpha + \delta) - F(\alpha) > \epsilon,
\]

und somit gibt es eine Lücke bei jedem $\alpha + \beta$, und $F$ müsste überall unstetig sein. Dies ist unmöglich, also muss $F$ stetig sein.
\end{proof}

\begin{lem}
  Suppose $F$ is a strictly increasing function defined on the nonnegative reals. Then $F$ is convex iff for all $\alpha, \beta, \gamma \geq 0$ the inequality from \ref{lemma-ungl} holds, i.e.,
$$
F(\alpha + \beta + \gamma) + F(\gamma) \geq F(\alpha + \gamma) + F(\beta + \gamma).
$$
\begin{proof}

$\Rightarrow:$ Angenommen, $F$ ist konvex, und sei $p = \alpha / (\alpha + \beta)$, dann gilt

\begin{align*}
p(\alpha + \beta + \gamma) + (1 - p)\gamma &= \alpha + \gamma\\
(1 - p)(\alpha + \beta + \gamma) + p\gamma &= \beta + \gamma.
\end{align*}

Damit gilt wegen Konvexität

\begin{align*}
F(\alpha + \gamma) + F(\beta + \gamma) &= F \left[ p (\alpha + \beta + \gamma) + (1 - p) \gamma \right] + F \left[ (1 - p) (\alpha + \beta + \gamma) + p \gamma \right] \\
&\leq p F(\alpha + \beta + \gamma) + (1 - p) F(\gamma) + (1 - p) F(\alpha + \beta + \gamma) + p F(\gamma) \\
&= F(\alpha + \beta + \gamma) + F(\gamma).
\end{align*}

$\Leftarrow:$
Um die Umkehrung zu beweisen, benutzen wir die Stetigkeit von $F$ aus Lemma \ref{lemma-stetig}. Angenommen, $\alpha > \beta$, dann

\begin{align*}
F(\alpha) + F(\beta) &= F \left( \frac{\alpha - \beta}{2} + \frac{\alpha - \beta}{2} + \beta \right) + F(\beta) \\
&\geq 2 F \left( \frac{\alpha - \beta}{2} + \beta \right) \quad \text{(nach Lemma 4)} \\
&= 2 F \left( \frac{\alpha + \beta}{2} \right),
\end{align*}

was äquivalent zur Konvexität für ein stetiges $F$ ist.
\end{proof}
\end{lem}

